% !TeX root = ../../Book2.tex
\section{差商方法}
\label{b2:sec:1601}

能量解只有一阶弱导数，不能直接对方程再求导。差商以两个平移值之差代替导数，先在已有的函数空间中操作，再由一致估计恢复二阶弱导数。这样既避免预设待证正则性，也清楚显示系数的变化会在估计中产生什么误差。

\subsection{差商}
\label{b2:subsec:160101}

设整数 \(d\geq1\)，\(\Omega\subset\mathbb R^d\) 为开集，\(u\in L^1_{\mathrm{loc}}(\Omega;\mathbb K)\)，其中 \(\mathbb K=\mathbb R\) 或 \(\mathbb C\)，\(k\in\{1,\ldots,d\}\)、\(h\in\mathbb R\setminus\{0\}\)。固定坐标方向 \(e_k\)，定义 \term[chashang]{差商}[difference quotient]
\[D_h^ku(x)=\frac{u(x+he_k)-u(x)}h\qquad(h\ne0).\]
只在所需平移仍留在定义域内时使用这一表达式。若 \(f,g\in C_c^\infty(\Omega)\)，并且它们支集的所需 \(h\) 平移都仍在 \(\Omega\) 内，换元给出离散分部积分；该式也适用于满足这些支撑条件的 \(L^2\) 函数
\[\int fD_{-h}^kg=-\int(D_h^kf)g.\]
其中 \(D_{-h}^kg=(g(x-he_k)-g(x))/(-h)\)；负号和负步长必须同时处理。乘积公式是
\[D_h^k(fg)=f(x+he_k)D_h^kg+(D_h^kf)g(x).\]

\Needspace{6\baselineskip}
\begin{lemma}[差商与弱导数]
\label{b2:lem:16-difference-quotient}
设整数 \(d\geq1\)，\(U\Subset V\subset\mathbb R^d\) 为非空开集，标量域 \(\mathbb K=\mathbb R\) 或 \(\mathbb C\)，\(k\in\{1,\ldots,d\}\)，\(h\in\mathbb R\setminus\{0\}\)，\(D_h^ku(x)=[u(x+he_k)-u(x)]/h\)，\(|h|<\operatorname{dist}(U,\partial V)\)（\(V=\mathbb R^d\) 时该距离取 \(+\infty\)）。若 \(u\in H^1(V;\mathbb K)\)，则
\[\lVert D_h^ku\rVert_{L^2(U)}\le\lVert\partial_ku\rVert_{L^2(V)}.\]
反之，若 \(u\in L^2(V;\mathbb K)\)，且存在 \(0\leq C<\infty\)、\(h_0>0\)，使 \(\lVert D_h^ku\rVert_{L^2(U)}\le C\) 对所有 \(0<|h|<h_0\) 且 \(|h|<\operatorname{dist}(U,\partial V)\) 成立，则 \(u\) 在 \(U\) 中有弱导数 \(\partial_ku\in L^2(U;\mathbb K)\)，其范数不超过 \(C\)。
\end{lemma}
\begin{proof}
对光滑 \(u\)，有 \(D_h^ku(x)=\int_0^1\partial_ku(x+the_k)\,\mathrm dt\)。Jensen 不等式、积分换序和各平移集合包含于 \(V\) 给出所列界。一般 \(H^1\) 函数先在包住所有相关线段的内部区域磨光，再以 \(H^1\) 收敛取极限，结论保持。

证明反向时，取 \(h_n\to0\)。\(L^2(U)\) 中的有界差商有弱收敛子列，记极限为 \(g\)。对 \(\varphi\in C_c^\infty(U)\)，离散分部积分和测试函数差商的一致收敛给出
\[\int_Ug\varphi=\lim_n\int_U(D_{h_n}^ku)\varphi
 =-\lim_n\int_VuD_{-h_n}^k\varphi=-\int_Uu\partial_k\varphi.\]
故 \(g\) 正是弱导数，弱下半连续性给出范数界。弱导数的唯一性说明所获对象不依赖抽取的子列。
\end{proof}

引理中的第一项是已有导数给出的上界，第二项则把统一差商界变成导数存在性；正则性证明需要把这两个方向配合使用。仅对某一个固定步长估计当然不能推出可微性。还应保留内部集合 \(U\Subset V\)：在边界附近任意平移会跑出区域，边界估计因此需要新的处理。

\subsection{内部 \texorpdfstring{\(H^2\)}{H2} 估计}
\label{b2:subsec:160102}

\Needspace{6\baselineskip}
\begin{theorem}[内部二阶估计]
\label{b2:thm:16-interior-h2}
设整数 \(d\geq1\)，\(\Omega\subset\mathbb R^d\) 为非空开集，\(r>0\)，\(B_{2r}=B(x_0,2r)\Subset\Omega\)，实矩阵场 \(A\in W^{1,\infty}(B_{2r};\mathbb R^{d\times d})\) 满足 \(\lVert A\rVert_\infty\le\Lambda<\infty\)，且对几乎处处 \(x\in B_{2r}\)、全部 \(\xi\in\mathbb R^d\) 有 \(\xi^{\mathsf T}A(x)\xi\ge\lambda|\xi|^2\)，其中 \(\lambda>0\)，矩阵取欧氏算子范数。若实值 \(u\in H^1(B_{2r};\mathbb R)\)、\(f\in L^2(B_{2r};\mathbb R)\) 且
\[\int_{B_{2r}}A\nabla u\cdot\nabla\varphi=\int_{B_{2r}}f\varphi
 \qquad(\varphi\in C_c^\infty(B_{2r};\mathbb R)),\]
则 \(u\in H^2(B_r)\)，并有
\[\lVert D^2u\rVert_{L^2(B_r)}
 \le C\left(\lVert f\rVert_{L^2(B_{2r})}
 +(r^{-1}+\lVert\nabla A\rVert_\infty)\cdot\lVert\nabla u\rVert_{L^2(B_{2r})}\right),\]
其中 \(C\) 只依赖 \(d,\lambda,\Lambda\)。
\end{theorem}
\begin{proof}
取光滑截断 \(0\le\eta\le1\)，使其在 \(B_r\) 上为 \(1\)，支集包含于 \(B_{3r/2}\)，且 \(|\nabla\eta|\le C_d/r\)。固定方向 \(k\)，取 \(|h|<r/8\)，记 \(w=D_h^ku\)。对这个固定步长，\(w\) 属于较小球上的 \(H^1\)，所以 \(v=-D_{-h}^k(\eta^2w)\) 属于 \(H_0^1(B_{2r})\)。型与右端均连续，密度保证可以用 \(v\) 测试。

将差商移到通量，得到
\[
 \int D_h^k(A\nabla u)\cdot\nabla(\eta^2w)
 =-\int fD_{-h}^k(\eta^2w).
\]
乘积公式给出 \(D_h^k(A\nabla u)=A(x+he_k)\nabla w+(D_h^kA)\nabla u\)。其中主项为 \(\int\eta^2A(x+he_k)\nabla w\cdot\nabla w\)，下界是 \(\lambda X^2\)，记 \(X=\lVert\eta\nabla w\rVert_2\)。设 \(U=\lVert\nabla u\rVert_{L^2(B_{2r})}\)、\(F=\lVert f\rVert_{L^2(B_{2r})}\)、\(K=\lVert\nabla A\rVert_\infty\)。引理给出 \(\lVert w\rVert_{L^2(B_{7r/4})}\le U\)，且 Lipschitz 控制给出 \(\lVert D_h^kA\rVert_\infty\le K\)。

展开 \(\nabla(\eta^2w)=\eta^2\nabla w+2\eta w\nabla\eta\)。其余主部交叉项、系数差商项及右端分别被
\[C\Lambda r^{-1}UX,\qquad KUX+CKr^{-1}U^2,
 \qquad CF(X+r^{-1}U)\]
控制。最后一界再次使用差商引理于紧支集函数 \(\eta^2w\)，即 \(\lVert D_{-h}^k(\eta^2w)\rVert_2\le\lVert\partial_k(\eta^2w)\rVert_2\)。逐项用 Young 不等式吸收不超过 \(\lambda X^2/2\) 的部分，余项得到
\[X^2\le C\bigl(F^2+(r^{-1}+K)^2U^2\bigr),\]
常数与 \(h\) 无关。因为 \(\eta=1\) 于 \(B_r\)，各 \(D_h^k\partial_j u\) 在该球的 \(L^2\) 范数一致有界。应用引理的反向得到全部 \(\partial_k\partial_j u\in L^2(B_r)\)，并对有限个指标求和即得估计。
\end{proof}

这个证明没有先假定二阶导数存在。测试函数中的差商仍只要求 \(u\in H^1\)；真正的新导数是在获得与步长无关的界之后才出现。系数的 Lipschitz 范数也有明确来源：它控制 \(D_hA\)。若系数只有本性上界，这个误差可能按 \(|h|^{-1}\) 增大，证明就不能闭合。

这里的矩阵系数取实值，解和右端的复值版本可直接分开处理：若 \(u=u_1+\mathrm i u_2\)、\(f=f_1+\mathrm i f_2\)，同一实矩阵 \(A\) 使两个实方程分别成立。各用一次定理，再把两个平方估计相加，就得到同一形式、仍只依赖 \(d,\lambda,\Lambda\) 的复值估计。下节的高阶提升也可逐次按这个办法处理复值数据；一般复矩阵的椭圆条件与配对须另作规定。

\subsection{局部正则性}
\label{b2:subsec:160103}

设整数 \(d\geq1\)，\(\Omega\subset\mathbb R^d\) 为非空开集，\(A\in W_{\mathrm{loc}}^{1,\infty}(\Omega;\mathbb R^{d\times d})\)，实值 \(f\in L_{\mathrm{loc}}^2(\Omega)\)、\(u\in H_{\mathrm{loc}}^1(\Omega)\) 满足 \(-\operatorname{div}(A\nabla u)=f\) 的局部测试恒等式。给定非空开集 \(U\Subset V\Subset\Omega\)，假设对几乎处处 \(x\in V\)、全部 \(\xi\in\mathbb R^d\) 有 \(\xi^{\mathsf T}A(x)\xi\geq\lambda_V|\xi|^2\)，其中 \(\lambda_V>0\)。用等半径内部球覆盖 \(\overline U\)，使加倍球均在 \(V\) 内、重叠次数只由 \(d\) 控制，再将球上估计平方求和，得到
\[\lVert u\rVert_{H^2(U)}\le C_{U,V,A}\bigl(\lVert f\rVert_{L^2(V)}+\lVert u\rVert_{H^1(V)}\bigr).\]
常数可只按 \(d\)、\(\operatorname{dist}(U,\partial V)^{-1}\)、\(\lambda_V\) 及 \(\|A\|_{W^{1,\infty}(V)}\) 选取。若每个紧含子域都有这样的椭圆下界，则局部 Lipschitz 系数和局部平方可积右端把能量解提高到 \(H_{\mathrm{loc}}^2\)。常数在区域逼近边界时可能发散，不能由此直接得到全域 \(H^2\) 估计。

\ProseParagraphBreak

若方程还含 \(b\cdot\nabla u+cu\)，且低阶系数有界，则把它移到右端：\(\widetilde f=f-b\cdot\nabla u-cu\in L_{\mathrm{loc}}^2\)。内部二阶估计仍适用，右端多出已有的 \(H^1\) 范数。这一步不需要对 \(b,c\) 求导；只有继续提高阶数时才会用到它们更高的正则性。

\ProseParagraphBreak

一个对照例是前章的分段导热系数：弱解通量连续而梯度跳跃，故二阶分布导数含点质量，不在 \(L^2\)。它说明系数正则性要求具有真实数学内容。另一方面，对常系数 Laplace 方程，本定理不要求 \(f\) 连续，只要求 \(f\in L^2\)；所得结论是二阶弱导数平方可积，而不是二阶经典导数连续。正则性结论应始终与数据所在的空间一起读取。
